% Generated draft: requires mathematical and source review.
\documentclass[11pt]{article}
\usepackage[margin=25mm]{geometry}
\usepackage{fontspec}
\setmainfont{Latin Modern Roman}
\usepackage{amsmath,amssymb,mathtools,unicode-math}
\setmathfont{Latin Modern Math}
\usepackage{enumitem,needspace,xurl,hyperref}
\setlength{\parindent}{0pt}
\setlength{\parskip}{5pt}
\setlength{\emergencystretch}{3em}
\newcommand{\N}{\mathbb{N}}
\newcommand{\Z}{\mathbb{Z}}
\newcommand{\Q}{\mathbb{Q}}
\newcommand{\R}{\mathbb{R}}
\newcommand{\C}{\mathbb{C}}
\newcommand{\F}{\mathbb{F}}
\newcommand{\A}{\mathbb{A}}
\newcommand{\PP}{\mathbb{P}}
\newcommand{\E}{\mathbb{E}}
\newcommand{\eps}{\varepsilon}
\newcommand{\dd}{\,\mathrm{d}}
\begin{document}
\section*{UnsolvedMath ID: 30005638}
\textbf{Source problem code:} catalogue entry
\subsection{Definitions and mathematical statement}
Fix an algebraic closure \(\overline{\Q}\) of \(\Q\). For a prime \(p\), let \(\mu_p\subset\overline{\Q}^{\times}\) be the group of \(p\)-th roots of unity, and put \(K_p=\Q(\mu_p)\). Define
\[
M_p=\bigcup_{L}L,
\]
where the union ranges over finite Galois extensions \(L/K_p\) inside \(\overline{\Q}\) whose Galois groups have \(p\)-power order and which are unramified at every finite prime of \(K_p\) not lying above \(p\). Thus \(M_p\) is the maximal pro-\(p\) extension with this finite-prime ramification restriction. The source does not specify its convention concerning archimedean ramification, which may matter when \(p=2\), or the intended range of primes.

\textbf{Historical question, incompletely specified in the supplied record.} With \(\Omega^{*}\) denoting the field intended by the source, is
\[
\Omega^{*}=M_p?
\]
The supplied excerpt does not define \(\Omega^{*}\), its ambient field, or its dependence on \(p\). Consequently this display records the historical target but is not a complete standalone conjecture. In particular, no definition of \(\Omega^{*}\) as a fixed field of a representation kernel can be inferred from the excerpt alone. The base field in the question is explicitly \(\Q(\mu_p)\).

\textbf{Associated structure reported by the source.} Write \(G_{\Q}=\operatorname{Gal}(\overline{\Q}/\Q)\). The excerpt refers to an outer representation \(\rho_{0,3}:G_{\Q}\to\operatorname{Out}(\Pi_{0,3})\) and a descending central series \(\Pi_{0,3}(m)\), and gives
\[
F^mG_{\Q}=\ker\!\left(G_{\Q}\xrightarrow{\rho_{0,3}}\operatorname{Out}(\Pi_{0,3})\longrightarrow\operatorname{Out}\bigl(\Pi_{0,3}/\Pi_{0,3}(m+1)\bigr)\right),\qquad m\geq1,
\]
\[
\operatorname{gr}^mG_{\Q}=F^mG_{\Q}/F^{m+1}G_{\Q},\qquad
\mathfrak g=\bigoplus_{m\geq1}\operatorname{gr}^mG_{\Q}.
\]
Here \(\operatorname{Out}\) denotes the automorphism group modulo inner automorphisms. The excerpt does not supply the construction of \(\Pi_{0,3}\), the representation, or their conventions. It reports that each graded piece is isomorphic to a finite direct sum of copies of the Tate twist \(\Z_p(m)\), and that commutators induce a graded Lie algebra structure on \(\mathfrak g\). Here \(\Z_p\) denotes the ring of \(p\)-adic integers, and \(\Z_p(m)\) its \(m\)-th cyclotomic Tate twist.

Put \(\mu_{p^{\infty}}=\bigcup_{n\geq1}\mu_{p^n}\) and \(G_{\Q(\mu_{p^{\infty}})}=\operatorname{Gal}(\overline{\Q}/\Q(\mu_{p^{\infty}}))\). For each odd integer \(m\geq3\), the excerpt reports a Soul\'e character
\[
\kappa_m:G_{\Q(\mu_{p^{\infty}})}\longrightarrow\Z_p(m)
\]
constructed as a Kummer character from cyclotomic \(p\)-units. It states that its restriction to \(F^mG_{\Q}\) is nontrivial and factors through \(\operatorname{gr}^mG_{\Q}\), and chooses \(\sigma_m\in\operatorname{gr}^mG_{\Q}\) such that \(\kappa_m(\sigma_m)\) generates \(\kappa_m(F^mG_{\Q})\); the excerpt calls this image an \(m\)-th Soul\'e element. These are reported background assertions, not independently verified statements in this draft.

\textbf{Unresolved source qualifications.} The excerpt says that the historical question has an affirmative answer under the Deligne--Ihara conjecture, described there as now a theorem of Hain--Matsumoto and Brown concerning an associated graded Lie algebra over \(\Q_p\). The supplied verification field repeats this implication. In contrast, the supplied dated literature assessment labels the problem partially solved and says that the general equality remains open, while citing Sharifi's theorem for odd irregular primes and the Deligne--Ihara implication in the relevant setting. The missing hypotheses and definition of \(\Omega^{*}\) prevent reconciliation of these statements. The excerpt also mentions a Deligne--Ihara analogue and an analogue of Sharifi's result for once-punctured CM elliptic curves over imaginary quadratic fields when \(p\) is ordinary; it gives no precise formulation of these analogues.
\subsection{Short English statement}
The historical Anderson--Ihara question asks whether the source's field \(\Omega^{*}\) equals the maximal pro-\(p\) extension of \(\Q(\mu_p)\) unramified outside \(p\). The supplied record omits the definition of \(\Omega^{*}\) and the admissible primes, and contains unresolved status qualifications. This is an incomplete statement-only draft requiring human mathematical review.
\subsection{Sources}
\begin{itemize}
\item Supplied local UnsolvedMath registry record, permanent ID \(30005638\), exact source problem code \texttt{catalogue entry}, title: ``Anderson-Ihara Maximal Pro-p Extension Conjecture.'' No proposer is supplied. The registry gives proposed year \(2023\), status \texttt{partially\_solved}, and statement status \texttt{unrecoverable}; these are supplied classifications, not independently verified mathematical conclusions.
\item Supplied report citation: \emph{MFO--RIMS Tandem Workshop: Arithmetic Homotopy and Galois Theory (2024)}, Oberwolfach Reports, pp.\ 2377--2488, DOI \texttt{10.4171/owr/2023/42}. \url{https://doi.org/10.4171/owr/2023/42}. This is also the supplied literature source labelled ``OWR report recording the known cases.'' The citation's year \(2024\) differs from the year component \(2023\) in the DOI and the registry's proposed year; this draft does not resolve that discrepancy.
\item Supplied literature source labelled ``Anderson--Ihara, pro-l branched coverings'': \url{https://annals.math.princeton.edu/1988/128-2/p05}. Full bibliographic details and a precise statement identifying \(\Omega^{*}\) are not included in the record.
\item The supplied excerpt cites Ichimura--Sakaguchi [IS87] for properties of Soul\'e characters, Rasmussen and Tamagawa [RT08] for further work, and Matsumoto, Nakamura, Takao and others for work related to Oda's problem. It refers to Philip's report in the same volume. These references are incomplete in the supplied record. Its attribution of the former Deligne--Ihara conjecture to a theorem of Hain--Matsumoto and Brown, and its mention of Sharifi's result, are retained as source attributions without added bibliographic details.
\item Literature checked at: \texttt{2026-08-22}. This date and the assessment are supplied dated curation, not independently verified. The assessment reports important known cases, including Sharifi's theorem for odd irregular primes, a Deligne--Ihara implication in the relevant setting, and an open general equality. The background explicitly requests independent verification of time-sensitive status and cited claims. The supplied statement-verification field additionally says that the report presents an affirmative answer under the former Deligne--Ihara conjecture, now a theorem, and that no separate clean statement was promoted while the original record was retained for provenance.
\end{itemize}
\end{document}
